% ============================================================
%   §1 Introduction
% ============================================================
\section{Introduction}
\label{sec:intro}

High-resolution spatial transcriptomics platforms such as 10x Genomics Visium HD now resolve gene expression at near single-cell resolution while preserving spatial coordinates and an aligned H\&E histological image \cite{visium_hd}. The resulting datasets are intrinsically network-shaped: every cell has neighbours, and those neighbours define a microenvironment that conditions its biology. Treating the data as a network --- with cells as nodes and spatial proximity as edges --- is not a methodological flourish but the natural framing of the object, and the lens through which I conduct every analysis in this report.

The dataset I work with comprises 8{,}065{,}561 individually segmented cells from a colorectal-cancer cohort of 56 patients (two samples per patient, 112 samples in total). Each cell carries a 16{,}497-dimensional gene expression vector, a 1280-dimensional H\&E image embedding produced by the UNI foundation model \cite{uni2024}, four morphological scalar features, and a categorical label drawn from 8 cell-type classes that span the tumour microenvironment: B cells (\classid{0}), Endothelial cells (\classid{1}), Fibroblasts (\classid{2}), Myeloid cells (\classid{3}), Normal epithelium (\classid{4}), SMC (\classid{5}), T cells (\classid{6}), and Tumor cells (\classid{7}). The cohort exhibits a 7.9$\times$ class imbalance, with Tumor cells and B cells over-represented and T cells and Endothelial cells under-represented. Spatial adjacency is encoded as a $k$-nearest-neighbour graph with $k = 8$, capped by a per-sample radius threshold. Figure~\ref{fig:intro_data_to_graph} illustrates the data-to-graph mapping on a representative sample: a tissue section becomes a point cloud coloured by cell type (panel a), zoomed regions reveal local cellular organisation (b), and the close-up exposes the KNN edges that connect each cell to its eight nearest neighbours (c).

\begin{figure}[h]
    \centering
    \includegraphics[width=0.97\linewidth]{intro_data_to_graph.png}
    \caption{From spatial transcriptomics data to a graph. (a) Sample 47 with all 50{,}988 cells coloured by cell type; the black rectangle marks the region zoomed in (b). (b) A 9{,}563-cell region of the same sample; the inner rectangle marks the close-up shown in (c). (c) A 715-cell close-up with 3{,}012 KNN edges drawn as grey lines. Edges follow spatial proximity rather than cell-type identity: a cell is connected to its eight nearest neighbours regardless of class.}
    \label{fig:intro_data_to_graph}
\end{figure}

I organise the project into five Parts, each posing a specific question on the cohort. \partref{1} asks what the spatial graph topology captures across the 112 samples: which metrics are consistent, which separate samples, which reflect biology rather than artefacts. \partref{2} asks whether the tissue organises into recurrent microenvironments by detecting Leiden communities in each sample and aligning them across patients into a stable set of niches. \partref{3} asks whether the spatial graph improves cell-type classification beyond per-cell features, evaluated through a five-rung ablation that progressively adds modalities (HVG expression, H\&E embedding, morphology) and architectures (logistic regression, MLP, GraphSAGE, GAT). \partref{4} asks whether each cell type carries its own gene co-expression network by fitting Gaussian Graphical Models on 378 broadly variable genes shared across classes. \partref{5} repeats the same fitting per-niche rather than per-class, then asks whether the resulting networks recover the composition-based niche families from \partref{2}.

Four findings stand out from the analysis. First, local graph topology is remarkably consistent across the cohort (mean degree near 8.2 for most samples), while a global metric --- the fraction of cells in the largest connected component --- splits the samples into two distinct modes, identifying compact and fragmented tissue architectures. Second, twenty niches recur stably across the cohort: 4 are universal (present in at least 85\% of patients), 11 common, 4 selective, and only one (a pure-SMC niche) is rare, appearing in just 11 patients. The 20 niches partition asymmetrically across the eight cell types: Tumor cells, B cells, and Fibroblasts each host four niches, while T cells host none. Third, the spatial graph improves validation macro-$F_1$ by 6.5 percentage points over a strong per-cell baseline (0.6350 vs.\ 0.5704), with a final test macro-$F_1$ of 0.6089, while adding the H\&E embedding, morphology features, or graph attention on top of GraphSAGE yields no further gain. Fourth, T cells emerge as the structurally distinctive class on every measure (topology, niche assignment, classification, gene network), with a unifying biological explanation: T cells in tumour tissue are diffuse infiltrating populations with no consistent spatial neighbourhood. Their one exception is \nicheid{17}, a B-cell-dominant niche where T cells coexist with B cells at substantial proportions --- a composition consistent with tertiary lymphoid structures, the only place in the niche landscape where T cells form an organised population. Across all 8 per-class and 20 per-niche gene networks, cell types and niches share dominant hubs (top-15 hub Jaccard up to 0.36 across classes, 0.58 across niches) but rewire the conditional-dependence connections between them (edge Jaccard at most 0.09 and 0.12 respectively).

The remainder of this report follows the five-Part structure (\S\ref{sec:part1}--\S\ref{sec:part5}), with a synthesis discussion in \S\ref{sec:discussion} and full methodological details in \S\ref{sec:methods}. Throughout, I refer to cell types by their biological labels (B cells, T cells, Tumor, etc.) in the discursive text, and add the numeric identifier (e.g.\ \classid{6}) at the first occurrence in each section and in cross-references to figures, which retain numeric \texttt{c0}--\texttt{c7} labels for consistency with the cached analysis. Niches are referred to as \nicheid{0}--\nicheid{19} with the dominant cell-type class annotated in parentheses in figures. Samples are indexed 0--111 and patients 0--55; the held-out test split is patients $\{0, 21, 26, 28, 34, 40, 45, 51\}$, fixed before any model training. One methodological qualification, important for reading Parts \ref{sec:part4}--\ref{sec:part5}: at the near single-cell resolution of Visium HD, mRNA from a target cell and its immediate spatial neighbours is difficult to disambiguate, and per-cell gene expression vectors carry a non-trivial component of neighbourhood signal. The implications of this diffusion contamination for the gene-network analyses are discussed in detail in \partref{4} and in \S\ref{sec:discussion}.